\subsection{Perturbation of Fredholm mappings}

THEOREM 3. --- Let E be a locally convex space, F a separated locally convex space, u a Fredholm map from E into F, and h a compact linear map from E into F. Then \(u + h\) is a Fredholm operator, and one has \(\operatorname{ind}(u + h) = \operatorname{ind}(u)\) .

We will prove the theorem in three steps.

\begin{enumerate}
\def\labelenumi{\Alph{enumi})}
\item
  Assume that E and F are Banach spaces. Prop. 2 of III, p. 42 shows that the linear map u is a strict morphism with closed image whose kernel and cokernel are finite-dimensional. Since E and F are Banach spaces, Ths. 1 and 2 of III, p. 73 imply that, for every \(t \in [0,1]\) , the linear map \(u_t = u + th\) possesses these same properties, and is therefore a Fredholm operator (III, p. 42, prop. 2). The map \(t \mapsto u_t\) from \([0,1]\) to the set \(\mathcal{F}(\mathrm{E};\mathrm{F})\) of Fredholm maps from E into F is continuous. By th. 1 of III, p. 58, the map \(t \mapsto \operatorname{ind}(u_t)\) is locally constant. Since the interval \([0,1]\) is connected, this map is constant. We therefore have \(\operatorname{ind}(u) = \operatorname{ind}(u + h)\) .
\item
  We assume that E = F. \(u = 1_{E}\) . By prop. 5 of III, p. 5, there exists a Banach space G, a continuous linear map \(p: E \to G\) and a compact linear map \(q: G \to E\) such that \(h = q \circ p\) . The endomorphism \(p \circ q\) of G is compact. By A), \(1_{G} + p \circ q\) is a Fredholm endomorphism of G with index zero. But then \(1_{E} + h = 1_{E} + q \circ p\) is a Fredholm endomorphism of E with index zero (III, p. 49, prop. 10, f)).
\item
  General case. Let \(v\) be a quasi-inverse of \(u\) . The endomorphisms \(u \circ v\) and \((u + h) \circ v\) of \(F\) differ from \(1_{F}\) by compact linear maps. By B), these are Fredholm endomorphisms of \(F\) of index zero. It follows that \(u + h\) is a Fredholm endomorphism (III, p. 40, n°2) with the same index as \(u\) (III, p. 43, no. 3, formula (2)).
\end{enumerate}

COROLLARY 1. --- Let E and F be separated locally convex spaces and \(u \in \mathcal{L}(\mathrm{E};\mathrm{F})\) . For u to be a Fredholm operator, it is necessary and sufficient that there exist a map \(v \in \mathcal{L}(\mathrm{F};\mathrm{E})\) such that the linear maps \(1_{\mathrm{E}} - v \circ u\) and \(1_{\mathrm{F}} - u \circ v\) are compact.

Since a continuous linear map of finite rank is compact, the stated condition is necessary.

Let \(v\) be an element of \(\mathcal{L}(\mathrm{F};\mathrm{E})\) such that the linear maps \(1_{\mathrm{E}} - v \circ u\) and \(1_{\mathrm{F}} - u \circ v\) are compact. By Theorem 3, \(v \circ u\) and \(u \circ v\) are Fredholm maps. Let \(p\) and \(q\) be quasi-inverses of \(v \circ u\) and \(u \circ v\) respectively. Set \(w = p \circ v\) and \(w' = v \circ q\) . Using the notation of No. 2 of III, p. 40, we have \(w \circ u \equiv 1_{\mathrm{E}}\) and \(u \circ w' \equiv 1_{\mathrm{F}}\) , whence \(w \equiv w \circ u \circ w' \equiv w'\) . It follows that \(w\) is a quasi-inverse of \(u\) .

COROLLARY 2. --- Let E and F be separated locally convex spaces and \(u \in \mathcal{L}(\mathrm{E};\mathrm{F})\) . The following conditions are equivalent:

\begin{enumerate}
\def\labelenumi{(\roman{enumi})}
\item
  The map u is a Fredholm map of index zero;
\item
  There exists an isomorphism v from E onto F such that u - v has finite rank;
\item
  There exists an isomorphism \(v\) of \(E\) onto \(F\) such that \(u - v\) is compact.
\item
  \(\Longrightarrow\) (ii): Suppose that u is a Fredholm operator of index zero. There exist topological direct sum decompositions \(E = E_{1} \oplus E_{2}\) and \(F = F_{1} \oplus F_{2}\) with \(E_{2}\) and \(F_{2}\) of finite dimension such that u vanishes on \(E_{2}\) and defines by restriction an isomorphism \(u_{1}\) of \(E_{1}\) onto \(F_{1}\) (III, p. 42, prop. 2). If \(\text{ind}(u) = 0\) then the dimensions of \(E_{2}\) and \(F_{2}\) are equal and there exists an isomorphism v from E onto F extending \(u_{1}\) . The kernel of u - v contains \(E_{1}\) , so u - v is of finite rank.
\end{enumerate}

We have (ii) \(\Longrightarrow\) (iii) since every continuous linear map of finite rank from E into F is compact, and (iii) \(\Longrightarrow\) (i) by Thm. 3.

